\documentclass[11pt,a4paper]{article}

\usepackage[a4paper,margin=1in]{geometry}
\usepackage{fontspec}
\IfFontExistsTF{Times New Roman}{\setmainfont{Times New Roman}}{\setmainfont{TeX Gyre Termes}}
\IfFontExistsTF{Menlo}{\setmonofont{Menlo}}{\setmonofont{Latin Modern Mono}}
\usepackage{amsmath,amssymb}
\usepackage{booktabs}
\usepackage{longtable}
\usepackage{array}
\usepackage{enumitem}
\usepackage{fancyhdr}
\usepackage[round,authoryear]{natbib}
\usepackage[hidelinks]{hyperref}
\usepackage{xurl}
\usepackage{microtype}

\setlength{\parindent}{0pt}
\setlength{\parskip}{0.65em}
\setlist[itemize]{leftmargin=1.4em}
\newcommand{\file}[1]{\path{#1}}
\newcommand{\code}[1]{\path{#1}}
\newcolumntype{L}[1]{>{\raggedright\arraybackslash}p{#1}}

\pagestyle{fancy}
\fancyhf{}
\lhead{\small Dataset Construction Plan}
\rhead{\small Planning Only}
\cfoot{\thepage}
\renewcommand{\headrulewidth}{0.4pt}

\title{\textbf{Dataset Construction Plan}\\
\large Final Planning Gate Before Any Dataset Generation}
\author{Codex-assisted surrogate planning report}
\date{June 18, 2026}

\begin{document}
\maketitle

\begin{abstract}
This report converts the approved surrogate dataset design into exact future construction
instructions. It approves construction planning only. It does not generate a dataset, solve the
planned regimes, create \code{.npz}, \code{.npy}, or \code{.pt} files, export labels, create
training tables, create model files, train neural networks, or run broad production-scale sweeps.
The future label source remains the validated baseline American Crank--Nicolson / projected SOR
solver.
\end{abstract}

\tableofcontents
\newpage

\section{Purpose of Dataset Construction Plan}

The surrogate dataset design plan approved only a schema and policy. This construction plan is
the next gate: it names the planned regimes, dry-run subset, split assignment, future storage
manifest, failure handling policy, and approval checklist. It is still a planning document. No
regime is solved here and no label array is produced.

The construction objective is a future dataset for constrained surrogate experiments on American
option value and continuation premium, with boundary and Greek diagnostics preserved as metadata
or masked labels. The project remains grounded in the validated finite-difference American option
benchmark and free-boundary setting
\citep{black_scholes_1973,merton_1973,brennan_schwartz_1977,wilmott_howison_dewynne_1995}.

\section{Design Plan Versus Construction Plan}

\begin{longtable}{L{0.28\textwidth} L{0.32\textwidth} L{0.28\textwidth}}
\toprule
Artifact & What it decides & What it does not do \\
\midrule
\endhead
Dataset design plan & Feature schema, label schema, parameter grid, masks, diagnostics, and split policy. & It does not name every regime or authorize generation. \\
Construction plan & Exact regime list, dry-run subset, split assignment, future storage manifest, and failure rules. & It does not run the solver or create dataset files. \\
Future dry run & A small future execution over eight regimes after approval. & It is not performed in this report. \\
Future small dataset & A later \code{v1_small_grid} generation if dry-run review passes. & It remains blocked until human review. \\
\bottomrule
\end{longtable}

\section{Source Evidence and Previous Gates}

The construction plan inherits constraints from the solver-validation synthesis, application
strength gate, downstream planning, Pilot 01, quantitative risk analytics, and the surrogate
dataset design plan. The methodology handouts warn that neural surrogates should be trained only
after the classical benchmark and application diagnostics are credible. They also warn against
random-only splits and unmasked Gamma labels near payoff kinks and free-boundary regions.

\begin{longtable}{L{0.30\textwidth} L{0.50\textwidth}}
\toprule
Prior source & Construction constraint \\
\midrule
\endhead
Ticket 12 synthesis & Solver validation passed with limitations, not production certification. \\
Application-strength gate & Dataset planning is allowed; dataset construction remains blocked. \\
Dataset design plan & Use normalized \(K=1\), baseline CN/PSOR, planned feature and label schemas, masks, diagnostics, and regime-level splits. \\
Methodology handouts & Use validated PSOR labels, stress regimes, held-out regimes, and boundary/Gamma caution. \\
\bottomrule
\end{longtable}

\section{Planned Full Regime Grid and Count}

The future full small-grid dataset is planned as the Cartesian product:
\[
\text{option\_type}\in\{\text{put},\text{call}\},\quad
T\in\{0.25,0.5,1.0,2.0\},\quad
\sigma\in\{0.20,0.40,0.60\},
\]
\[
r\in\{0.01,0.05,0.10\},\qquad
q\in\{0.00,0.03,0.06,0.10\}.
\]

This gives:
\[
2 \times 4 \times 3 \times 3 \times 4 = 288
\]
planned regimes: 144 American puts and 144 American calls. The 36 no-dividend call regimes are
included as validation/control regimes. They are not core dividend-call early-exercise evidence.

Every planned regime uses:
\begin{center}
\(K=1,\quad S_{\max}=4K,\quad M=N=120\),\\
\code{solver_variant = baseline_cn_psor}.
\end{center}

\begin{longtable}{L{0.27\textwidth} L{0.53\textwidth}}
\toprule
Regime-plan rule & Exact policy \\
\midrule
\endhead
Regime ID format & \code{\{option\}_T\{025|050|100|200\}_s\{020|040|060\}_r\{001|005|010\}_q\{000|003|006|010\}}. \\
Machine-readable plan & \file{reports/05_surrogate/dataset_regime_plan.csv}. \\
Generated now & No. The CSV is a plan, not solved data. \\
Control rows & Call rows with \(q=0\) are marked as no-dividend call controls. \\
\bottomrule
\end{longtable}

\section{Dry-Run Subset Plan}

Before any 288-regime generation, the future implementation must run a small manual-review dry
run. The dry run is planned as eight regimes:

\begin{longtable}{L{0.23\textwidth} L{0.31\textwidth} L{0.30\textwidth}}
\toprule
Dry-run ID & Regime ID & Purpose \\
\midrule
\endhead
\code{dry_01} & \code{put_T100_s020_r005_q003} & Base-like American put. \\
\code{dry_02} & \code{put_T100_s060_r005_q003} & High-volatility put. \\
\code{dry_03} & \code{put_T100_s020_r001_q003} & Low-rate put. \\
\code{dry_04} & \code{put_T200_s040_r010_q010} & Longer, higher-rate stress put. \\
\code{dry_05} & \code{call_T100_s020_r005_q006} & Base-like dividend-paying call. \\
\code{dry_06} & \code{call_T100_s020_r005_q010} & High-dividend call. \\
\code{dry_07} & \code{call_T100_s060_r005_q006} & High-volatility dividend call. \\
\code{dry_08} & \code{call_T100_s020_r005_q000} & No-dividend call control. \\
\bottomrule
\end{longtable}

The machine-readable dry-run plan is listed in the reproducibility and source record. It is not a
dry-run result table.

\section{Feature and Label Schema Summary}

The future feature vector remains:
\[
x=\left(\log(S/K),\ \tau/T,\ r,\ q,\ \sigma,\ T,\ \text{is\_call}\right),
\]
with solver and grid metadata retained separately. The main learning and reporting region remains:
\[
S/K\in[0.4,1.8],\qquad \tau/T\in[0,1].
\]

The most reliable future labels are \code{value_over_K} with formula \(U/K\),
\code{payoff_over_K} with formula \(\Phi/K\), and \code{premium_over_K} with formula
\((U-\Phi)/K\). Exercise indicators and boundary locations are threshold-based diagnostics.
Delta is cautioned. Scaled Gamma is high caution and must never be used unmasked near payoff
kinks, maturity rows, or extracted boundary bands.

\section{Mask and Diagnostic Policy Summary}

Future generated surfaces must preserve masks for payoff-kink-near, boundary-near, maturity-row,
strict-interior, Gamma-allowed, Delta-allowed, exercise-region, and continuation-region nodes.
Masks are mandatory metadata, not decoration.

Every future solved regime must record:
\begin{itemize}
  \item solver convergence flag;
  \item maximum and mean PSOR iterations;
  \item maximum final PSOR update;
  \item maximum obstacle violation;
  \item maximum equation violation;
  \item maximum absolute complementarity product;
  \item boundary found count and boundary threshold;
  \item grid metadata \(M,N,S_{\max},\Delta S,\Delta\tau\);
  \item solver name and solver variant;
  \item downstream use status.
\end{itemize}

\section{Split Assignment Plan}

The split is regime-level, not row-level and not random point only. The deterministic precedence is:
\begin{enumerate}
  \item \code{stress_holdout}: \(\sigma=0.60\) and \(q=0.10\);
  \item \code{test}: all call regimes with \(q=0.00\), plus \(T=2.0\), \(r=0.10\), \(\sigma=0.40\), after stress precedence;
  \item \code{validation}: \(T=0.5\) and \(r=0.05\), after stress/test precedence;
  \item \code{train}: all remaining regimes.
\end{enumerate}

\begin{longtable}{L{0.24\textwidth} L{0.20\textwidth} L{0.42\textwidth}}
\toprule
Split & Planned count & Purpose \\
\midrule
\endhead
\code{train} & 202 & Main fitting regimes for future surrogate experiments. \\
\code{validation} & 19 & Model-selection regimes after stress/test precedence. \\
\code{test} & 43 & No-dividend call controls and selected long-maturity high-rate regimes. \\
\code{stress_holdout} & 24 & High-volatility, high-dividend stress regimes. \\
\bottomrule
\end{longtable}

The machine-readable assignment is listed in the reproducibility and source record. No generated
rows exist yet.

\section{Storage and Manifest Plan}

The future storage package may use one compressed \code{.npz} numeric-array file plus CSV or JSON
manifests. This report does not create that storage. The planned storage manifest has
\code{created_now=no} for every row.

\begin{longtable}{L{0.28\textwidth} L{0.22\textwidth} L{0.34\textwidth}}
\toprule
Future file & Type & Purpose \\
\midrule
\endhead
Future arrays file & compressed arrays & Future numeric samples only after approval; exact \code{.npz} name is in the CSV plan. \\
Regime manifest & CSV manifest & Regime parameters and solver settings. \\
Diagnostic summary & CSV manifest & PSOR, LCP, boundary, runtime, and mask diagnostics. \\
Split assignment & CSV manifest & Approved regime-level split assignment. \\
Label schema copy & CSV schema copy & Frozen future label schema. \\
Mask policy & JSON manifest & Exact mask thresholds and widths. \\
Generation log & text log & Future command, environment, and failure log. \\
Validation report & PDF report & Future post-generation review report. \\
\bottomrule
\end{longtable}

\section{Failure Handling Policy}

The future generation script must reject or quarantine regimes according to the failure policy
before any label export:

\begin{longtable}{L{0.30\textwidth} L{0.34\textwidth} L{0.22\textwidth}}
\toprule
Failure type & Detection rule & Default action \\
\midrule
\endhead
PSOR non-convergence & PSOR convergence flag is false. & Exclude and review. \\
Obstacle violation & Obstacle violation exceeds \(10^{-8}\). & Exclude and review. \\
Equation violation & Equation violation exceeds \(10^{-6}\). & Exclude and review. \\
Complementarity failure & Complementarity product exceeds \(10^{-6}\). & Exclude and review. \\
Boundary not found & Boundary flag false or count unexpectedly low. & Keep price/premium only if other diagnostics pass; boundary label missing. \\
Greek masks missing & Required masks absent or unapproved. & Block Delta/Gamma labels. \\
Unexpected runtime & Runtime exceeds approved budget. & Mark expensive and review. \\
Missing metadata & Required provenance or diagnostics absent. & Block export until fixed. \\
\bottomrule
\end{longtable}

\section{Construction Approval Checklist}

Actual dataset construction may start only after human review approves:
\begin{itemize}
  \item exact regimes and 288-regime count;
  \item dry-run subset;
  \item split assignment;
  \item storage format;
  \item diagnostic thresholds;
  \item failure handling policy;
  \item mask widths;
  \item post-dry-run review process.
\end{itemize}

The checklist is listed in the reproducibility and source record. It records planning status only.

\section{Blocked Actions}

The following remain blocked by this plan:
\begin{itemize}
  \item actual dataset construction;
  \item CN/PSOR solves over the planned regimes;
  \item \code{.npz}, \code{.npy}, or \code{.pt} file creation;
  \item label-array export;
  \item training-table export;
  \item model-file creation;
  \item neural surrogate training;
  \item broad production-scale sweeps;
  \item production Greek labels;
  \item final paper claims across all regimes.
\end{itemize}

\section{Limitations}

This construction plan is detailed enough to implement a future dry-run generator, but it is not
evidence that the dataset exists. The split counts are deterministic and auditable, but they have
not been tested on generated rows. Boundary and Gamma labels remain diagnostic and mask-dependent.
The planned \(M=N=120\) grid is inherited from pilot defaults and should be confirmed on selected
\(M=N=180\) cases before broad scaling.

\section{Recommended Next Step}

The recommended next step is human review of this construction plan and its six planning CSVs. If
accepted, the next technical ticket should implement only the eight-regime \code{v0_dry_run}
generator, parse the approved planning CSVs, write no full 288-regime dataset, and produce a
post-dry-run review report before any \code{v1_small_grid} generation.

\section{Reproducibility and Source Record}

\begin{longtable}{L{0.34\textwidth} L{0.48\textwidth}}
\toprule
Item & Record \\
\midrule
\endhead
Regime plan & \file{reports/05_surrogate/dataset_regime_plan.csv}. \\
Dry-run subset & \file{reports/05_surrogate/dataset_dry_run_subset_plan.csv}. \\
Split assignment & \file{reports/05_surrogate/dataset_split_assignment_plan.csv}. \\
Storage manifest plan & \file{reports/05_surrogate/dataset_storage_manifest_plan.csv}. \\
Failure handling policy & \file{reports/05_surrogate/dataset_failure_handling_policy.csv}. \\
Construction checklist & \file{reports/05_surrogate/dataset_construction_checklist.csv}. \\
Construction backlog & \file{reports/05_surrogate/dataset_construction_backlog.md}. \\
Shared bibliography & \file{reports/03_solver/references.bib}. \\
\bottomrule
\end{longtable}

\bibliographystyle{plainnat}
\bibliography{reports/03_solver/references}

\end{document}
